\section{实验设计与结果分析}

\subsection{实验说明与可复现性}
本章的任务书对应关系可以概括为两条并行主线：\textbf{主线 A（影像）}——在 CTA 体数据上训练三维 Attention U-Net，输出与金标准可比对的分割概率图与二值掩膜；\textbf{主线 B（表型）}——在 Excel/CSV 表格上训练破裂风险二分类模型，输出校准前的风险概率及测试集评价摘要。两条主线在工程上通过可选的形态学/组学特征字段衔接，在答辩演示上通过 Web 界面串联。

需要强调的是：受医院数据脱敏、伦理审批进度与算力排期等限制，不同复现者在 Dice、AUC 等具体数值上可能存在差异；\textbf{正式归档时，应以经审批数据集上、由本项目脚本与日志导出的客观记录为准}。下文除特别标注外，凡出现“示例性”表格或百分比，均可在定稿前整体替换为实测；章节后半部分则尽量扩充\textbf{与代码一致的过程性描述}，避免仅有少量数字而缺少实验设计交代。

\subsection{实验环境与主要依赖}
\subsubsection{硬件与软件}
本课题开发与调试的典型环境如下（答辩材料中可按实验室机器实录微调）：
\begin{itemize}
	\item 操作系统：Linux（Ubuntu 20.04 LTS 或兼容发行版）；
	\item GPU：NVIDIA CUDA 兼容显卡（例：RTX 3090，24GB 显存），用于三维卷积与批量矩阵运算；
	\item Python：3.8 及以上（建议与 PyTorch 官方 wheel 版本说明一致）；
	\item 深度学习框架：PyTorch，配合 MONAI 中的体数据变换、度量实现（若实际未调用 MONAI API，也可在论文中写“计划对接 MONAI 指标以统一口径”）；
	\item 医学图像 IO：SimpleITK，用于读写 NIfTI、保留 spacing/origin/direction；
	\item 实验管理：TensorBoard（训练曲线）、\texttt{tqdm}（终端进度条）、YAML（超参）；
	\item 风险任务：Pandas、scikit-learn（划分、互信息、ROC/AUC）；
	\item 演示服务：FastAPI、Uvicorn。
\end{itemize}

设备由配置文件中的 \texttt{device} 指定，并可用命令行 \texttt{--device cuda|cpu} 覆盖，便于笔记本 CPU 调试与服务器 GPU 训练切换。

\subsubsection{随机种子与版本记录}
为便于问题复现，项目在 \texttt{core} 层对随机种子进行统一设置（如 NumPy、PyTorch CUDA）。建议在毕业论文附录中\textbf{固定给出}：Python 与 PyTorch 版本号、CUDA/cuDNN 版本、关键依赖包版本、所用配置文件路径（如 \texttt{config/remote.yaml}）及 Git 提交哈希（若适用）。答辩评委追问“结果能否复现”时，上述信息是最直接的技术答复。

\subsection{数据协议与预处理约定}
\subsubsection{分割任务数据}
动脉瘤分割数据以 NIfTI（\texttt{.nii} 或 \texttt{.nii.gz}）为主。训练与验证目录由 \texttt{data.train\_dirs}、\texttt{data.eval\_dirs} 分别列出；若未单独指定验证目录，代码侧可能回退到训练目录（日志中会出现警告），此时验证 Dice 的独立性较弱，论文中应如实说明划分策略。

图像与标签通过文件名规则配对：\texttt{file\_patterns} 列出影像候选模式（如 \texttt{*\_origin.nii.gz}），\texttt{label\_suffix} 列出标签后缀（如 \texttt{\_label.nii.gz}、\texttt{\_ias.nii.gz}）。该设计允许同一批原始数据以不同预处理版本共存，只要后缀规则可唯一推断标签路径。

体素各向异性通过重采样统一到目标 spacing；强度在补丁读取管线中标准化，以减轻不同扫描参数带来的灰度分布偏移。与二维分类数据集“固定张数图片”不同，体数据实验需额外记录：\textbf{纳入病例数}、\textbf{每例层厚/重建矩阵}、\textbf{标注金标准来源}（单专家/双专家融合）、\textbf{最小动脉瘤直径纳入标准}等，否则难以与其他文献横向对比。

\subsubsection{破裂风险任务数据}
风险建模依赖表格文件路径 \texttt{risk.excel\_path}、标签列 \texttt{label\_column}（如“破裂”）、病例 ID 列 \texttt{id\_column}。特征列可以显式枚举，也可以留空触发基于互信息（\texttt{feature\_select\_method: mutual\_info}）的自动筛选，并受 \texttt{auto\_max\_features} 上限约束。配置中还包含临床先验列（如性别、年龄、合并症等）、类别低频合并阈值（\texttt{categorical\_min\_frequency}）以及是否与影像组学特征对齐（\texttt{use\_pyradiomics\_overlap\_only}）等开关——这些字段决定了\textbf{实际输入模型的特征数量与语义}，论文中建议用表格列出“最终进入训练的数值列数、类别列数、各类基数”，而不是笼统写“数十个特征”。

\subsection{主线 A：动脉瘤分割实验设计}
\subsubsection{训练协议与超参数（与配置项对齐）}
分割训练入口为 \texttt{script/train\_split.py}，核心超参由 YAML 的 \texttt{model} 与 \texttt{train} 段驱动。表~\ref{tab:seg_hyper} 归纳了与本仓库默认思路一致、且应在论文中逐项交代的关键项（具体取值以你实际使用的配置文件为准）。

\begin{table}[htbp]
	\centering
	\caption{分割任务关键超参数与典型取值（请按实际 YAML 填写最终列）}
	\label{tab:seg_hyper}
	\begin{tabular}{p{4.2cm}p{7.8cm}}
		\toprule
		配置键 & 含义与实验设计说明 \\
		\midrule
		\texttt{model.depth}, \texttt{model.base\_filters} & 控制网络宽度与深度；深度越大感受野越大，但显存与过拟合风险同步上升。 \\
		\texttt{model.norm\_type}, \texttt{model.activation}, \texttt{model.dropout} & 归一化与激活选择影响小批量统计稳定性；Dropout 在体数据中需谨慎开启。 \\
		\texttt{train.epochs} & 总轮数；需结合 \texttt{eval\_interval} 观察验证 Dice 是否仍上升，避免无效延长训练。 \\
		\texttt{train.batch\_size} & 本项目训练循环按“病例”迭代，常见取值为 1；应在文中解释为何不以更大 batch 堆叠多病例。 \\
		\texttt{train.patch\_size} & 三维补丁边长，直接决定显存占用；可在消融中对比 $96^3$ 与 $128^3$ 等。 \\
		\texttt{train.max\_patches\_per\_volume}, \texttt{data.patches\_per\_volume} & 每例最多采样补丁数，影响单 epoch 时长与前景覆盖。 \\
		\texttt{data.patch\_sampling\_mode} & 训练用 \texttt{foreground\_priority} 可提高前景占比；验证/整例推理常改为 \texttt{sequential}。 \\
		\texttt{data.background\_per\_foreground} & 前景补丁与背景补丁配比，用于压制“全预测背景”的平凡解。 \\
		\texttt{train.optimizer}, \texttt{train.learning\_rate}, \texttt{train.weight\_decay} & 优化器选择与正则；可在附录给出学习率曲线截图。 \\
		\texttt{train.scheduler}, \texttt{train.warmup\_epochs} & 学习率调度策略；若代码中 scheduler 步进尚未完全接线，应在论文中如实说明当前实际生效部分。 \\
		\texttt{eval\_interval} & 每隔多少 epoch 做一次验证并刷新 \texttt{best.pth}。 \\
		\texttt{checkpoint.save\_interval} & 周期性保存 \texttt{latest.pth}，用于防止中断后回滚。 \\
		\bottomrule
	\end{tabular}
\end{table}

\subsubsection{训练过程监控与产出文件}
每轮训练将 epoch 损失写入 TensorBoard 标量（目录一般为 \texttt{log\_dir/tensorboard}）。验证阶段由 \texttt{eval\_split.evaluate} 计算 Dice，并触发 \texttt{save\_checkpoint}：始终更新 \texttt{checkpoints/latest.pth}；若当前验证 Dice 高于历史最优，则写入 \texttt{checkpoints/best.pth}。论文中可截取 TensorBoard 中 \texttt{Loss/train\_epoch} 与验证 Dice 曲线各一张，作为“训练收敛行为”的定性证据。

此外，脚本可选保存训练/验证过程中的补丁预测预览（路径与 \texttt{prediction\_dir} 相关），便于肉眼排查标签颠倒、强度窗口错误等低级问题。

\subsubsection{整例推理与空间一致性}
验证或离线测试时，数据集将补丁采样模式切换为顺序遍历，并对整例体积进行切块推理与拼接；随后利用 SimpleITK 写回 NIfTI，并恢复原始影像的 spacing、origin、direction。这一环节是\textbf{临床可用性}的关键：若缺失几何元数据，后续与放疗计划系统或手术导航叠加将失败。实验章节建议单独用一小节文字说明：推理输出包含\textbf{概率图}与\textbf{二值掩膜}两类文件，及其后缀命名约定。

\subsubsection{分割结果的多维度解读（建议在答辩中展开）}
即使暂不报告大量病例统计，也可以在文中预先搭建分析框架：
\begin{itemize}
	\item \textbf{整体层面}：病例平均 Dice、Dice 的中位数与四分位距，反映算法稳定性；
	\item \textbf{病灶尺度分层}：按瘤体等效直径将病例分桶，讨论小动脉瘤是否 Dice 偏低（部分体素过少导致Dice方差大）；
	\item \textbf{解剖部位分层}：Willis 环与前循环、后循环动脉瘤在血管曲率与对比剂充盈上的差异，可能导致分割难易不同；
	\item \textbf{错误类型}：假阴性（漏分割）与假阳性（血管锥或邻近骨质高密度误判）在临床上后果不同，可分别统计体素级敏感度/特异度；
	\item \textbf{边界指标}：除 DSC 外，报告 95\% Hausdorff 距离，敏感于瘤颈区域的细小偏差。
\end{itemize}

\subsubsection{横向对比与消融实验框架}
表~\ref{tab:seg_result} 给出分割任务常见的横向对比排版方式。\textbf{基线模型行}可通过三种途径之一填充：（1）在同一数据划分上复现文献公开配置；（2）在项目中临时替换 \texttt{build\_model} 为纯 3D U-Net 并保持相同补丁协议；（3）若工期不足，则删除基线数值，仅保留本文方法与定性讨论。

\begin{table}[htbp]
	\centering
	\caption{动脉瘤分割任务：方法—指标对比表（示例性数值，提交前请替换为实测）}
	\label{tab:seg_result}
	\begin{tabular}{cccccc}
		\toprule
		方法 & DSC(\%) & 95\%HD(mm) & 敏感性(\%) & 特异性(\%) & 训练时间(h)\\
		\midrule
		3D U-Net & 78.3 & 4.2 & 76.8 & 94.2 & 12.5 \\
		3D U-Net++ & 81.2 & 3.6 & 79.5 & 94.8 & 16.8 \\
		V-Net & 79.6 & 3.9 & 78.2 & 94.5 & 14.2 \\
		本文 Attention U-Net & \textbf{85.7} & \textbf{2.8} & \textbf{84.1} & \textbf{96.3} & 15.3 \\
		\bottomrule
	\end{tabular}
\end{table}

消融实验建议覆盖第二章涉及的机制：注意力门控是否启用、损失是否包含 Dice、补丁采样是否前景优先、是否使用深度监督（若代码未实现则不予编造）。表~\ref{tab:ablation_result} 给出记录模板。

\subsection{主线 B：破裂风险预测实验设计}
\subsubsection{数据划分与泄漏防护}
\texttt{train\_risk.py} 通过 \texttt{test\_size}、\texttt{val\_size} 与 \texttt{random\_state} 完成训练集、验证集与测试集的划分。论文中应明确：互信息特征筛选\textbf{仅在训练集上拟合}，再将同一变换应用到验证与测试；若误用全部样本做筛选，会造成标签泄漏，AUC 虚高。若课题使用病例 ID 去重或家族簇抽样，也需在文中说明，否则表格任务的“独立同分布”假设不成立。

\subsubsection{模型超参数与训练细节}
表~\ref{tab:risk_hyper} 列出风险模块在配置中的主要接口，与 \texttt{RiskCrossAttentionModel} 构造函数一一对应。

\begin{table}[htbp]
	\centering
	\caption{破裂风险任务关键超参数与含义（按实际 YAML 填写取值）}
	\label{tab:risk_hyper}
	\begin{tabular}{p{4cm}p{8cm}}
		\toprule
		配置键 & 含义 \\
		\midrule
		\texttt{risk.hidden\_dim} & 隐空间维度，需与 \texttt{num\_heads} 整除约束一致。 \\
		\texttt{risk.num\_heads} & 多头注意力头数；头数过多在小样本下可能过拟合。 \\
		\texttt{risk.num\_layers} & 双向交叉注意力堆叠层数，depth 越大模型容量越大。 \\
		\texttt{risk.dropout} & Dropout 比率，配合 LayerNorm 与 GELU 缓解过拟合。 \\
		\texttt{risk.epochs}, \texttt{risk.batch\_size} & 表格任务通常可采用较大 batch；若类别极不均衡，可适当减小 batch 以增加梯度噪声、配合加权损失。 \\
		\texttt{risk.learning\_rate}, \texttt{risk.weight\_decay} & 优化器通常采用 AdamW；权重衰减对嵌入表同样生效。 \\
		\texttt{risk.use\_pos\_weight} & 是否根据训练集正负比自动设置正类权重。 \\
		\texttt{risk.threshold\_search\_steps} & 验证集上阈值扫描的离散点数；影响 F1 最优解的分辨率。 \\
		\texttt{risk.early\_stopping\_patience} 等 & 早停策略，防止验证集指标长期不升时过训练。 \\
		\bottomrule
	\end{tabular}
\end{table}

\subsubsection{优化目标与决策阈值}
训练损失为加权二元交叉熵，正类权重由训练集样本频次估计。验证阶段在每个 epoch 结束后，将验证集预测概率在 $(0.05,0.95)$ 上做多阈值搜索，以\textbf{宏平均 F1} 为主准则更新当前阈值；测试阶段固定验证集上选出的最优阈值，输出准确率、精确率、召回率与 AUC。论文中建议补充一句：若临床更重视“不漏诊破裂风险”，可在讨论中改用召回率或加权代价敏感准则重新搜索阈值。

\subsubsection{产出物与结果登记表}
训练结束后，目录 \texttt{risk.save\_dir} 下通常生成：
\begin{itemize}
	\item \texttt{risk\_cross\_attention.pt}：模型权重与还原所需的列名、类别基数等元数据；
	\item \texttt{risk\_report.json}：验证集最优 F1、AUC、阈值及测试集指标摘要。
\end{itemize}
答辩材料可将 JSON 中字段摘抄为表~\ref{tab:risk_result} 的“本文方法”行，避免手工抄错。

\begin{table}[htbp]
	\centering
	\caption{破裂风险预测任务：方法—指标对比表（示例性数值，提交前请替换为 risk\_report.json 实测）}
	\label{tab:risk_result}
	\begin{tabular}{ccccccc}
		\toprule
		方法 & 准确率(\%) & 精确率(\%) & 召回率(\%) & F1分数(\%) & AUC & 训练时间(min)\\
		\midrule
		逻辑回归 & 72.4 & 69.8 & 73.2 & 71.5 & 0.756 & 2.3 \\
		SVM & 73.8 & 71.2 & 72.6 & 71.9 & 0.763 & 4.7 \\
		Random Forest & 74.6 & 72.1 & 74.3 & 73.2 & 0.771 & 3.8 \\
		XGBoost & 76.2 & 73.5 & 75.8 & 74.6 & 0.785 & 5.2 \\
		本文交叉注意力模型 & \textbf{81.4} & \textbf{79.3} & \textbf{82.1} & \textbf{80.7} & \textbf{0.842} & 8.7 \\
		\bottomrule
	\end{tabular}
\end{table}

\subsubsection{风险模型的结果解读要点}
除整体 AUC 外，建议在讨论中预留以下分析维度：
\begin{itemize}
	\item \textbf{校准性}：预测概率是否系统性偏高/偏低（可通过可靠性曲线或分箱后的 observed rate 评估）；
	\item \textbf{少数类表现}：破裂阳性样本较少时，单独报告阳性子集的敏感度；
	\item \textbf{特征维度敏感性}：改变 \texttt{auto\_max\_features} 或手动剔除某类临床字段后，AUC 变化是否剧烈，用以判断模型是否依赖个别字段；
	\item \textbf{与分割衔接}：若 UI 或后处理脚本将分割体积导入形态学特征（如瘤体体积、长径），应说明该特征是否纳入风险模型及噪声来源（分割误差传递）。
\end{itemize}

\subsection{系统集成与演示实验}
\subsubsection{Web 接口层面的任务闭环}
\texttt{ui/main.py} 提供 HTTP API：上传 CTA 体积后调用与训练一致的补丁推理逻辑；风险预测端加载 \texttt{risk\_cross\_attention.pt} 并对表单字段编码。论文中可增加一小段“演示实验”：记录一次端到端请求耗时、占用显存峰值、返回 JSON 字段是否齐全。该类实验不替代离线批量评测，但能证明\textbf{集成链路可用}。

\subsubsection{任务书交付物对照表}
为方便评委核对“是否完成题目要求”，可将交付物梳理为表~\ref{tab:deliverables}（路径按仓库实际为准）。

\begin{table}[htbp]
	\centering
	\caption{毕业设计任务与仓库产出物对照（示例）}
	\label{tab:deliverables}
	\begin{tabular}{p{3.6cm}p{9cm}}
		\toprule
		任务书要求 & 对应实现与佐证材料 \\
		\midrule
		动脉瘤分割/检测 & \texttt{train\_split.py}、\texttt{eval\_split.py}、\texttt{best.pth}、验证 Dice 曲线、NIfTI 预测结果 \\
		破裂风险预测 & \texttt{train\_risk.py}（需 \texttt{risk.enabled: true}）、\texttt{risk\_report.json}、混淆矩阵（可自行脚本绘制） \\
		可视化或交互 & FastAPI UI、上传接口日志、浏览器截图 \\
		可复现实验 & YAML 配置、随机种子、依赖版本说明、附录训练命令 \\
		\bottomrule
	\end{tabular}
\end{table}

\subsection{消融实验与灵敏度分析}
表~\ref{tab:ablation_result} 汇总分割与风险两侧可做的消融维度。实施原则是：每次只改一个因素，固定随机种子与数据划分。

\begin{table}[htbp]
	\centering
	\caption{消融实验记录模板（请替换为实测）}
	\label{tab:ablation_result}
	\begin{tabular}{ccccc}
		\toprule
		配置 & DSC(\%) & 95\%HD(mm) & 准确率(\%) & AUC \\
		\midrule
		去除注意力门控 & 78.3 & 4.2 & - & - \\
		仅 BCE、无 Dice & 83.1 & 3.2 & - & - \\
		完整分割管线 & \textbf{85.7} & \textbf{2.8} & - & - \\
		\midrule
		浅层 MLP 风险基线 & - & - & 76.2 & 0.785 \\
		加入交叉注意力层 & - & - & 79.8 & 0.821 \\
		再加入验证集阈值搜索 & - & - & \textbf{81.4} & \textbf{0.842} \\
		\bottomrule
	\end{tabular}
\end{table}

灵敏度分析建议至少覆盖：\texttt{patch\_size}、\texttt{patches\_per\_volume}、\texttt{risk.hidden\_dim}、互信息保留特征数。可将横轴为超参取值、纵轴为验证指标的折线图附在答辩附录。

\subsection{可视化与定性案例分析}
图~\ref{fig:case_study} 建议展示至少两类病例：\textbf{Dice 较高}的代表性成功病例，与\textbf{Dice 偏低或存在临床关切误差}的困难病例；每张图旁配以简短文字说明可能原因（瘤体过小、运动伪影、邻近骨质干扰等）。

\begin{figure}[htbp]
	\centering
	\fbox{\parbox{0.72\textwidth}{\centering\vspace{2.8cm}\songti 此处插入分割可视化效果图（最大密度投影或轴位/冠状位拼图；建议标注金标准轮廓与预测轮廓）\vspace{2.8cm}}}
	\caption{典型分割结果示意（占位：可替换为 \texttt{includegraphics} 指向实测截图）}
	\label{fig:case_study}
\end{figure}

\subsection{本章小结}
本章在保留“评价指标 + 对比实验表”的主干之外，补充了与\textbf{实际代码一致}的训练协议、配置键说明、监控产物路径、风险划分与泄漏防护、阈值搜索逻辑、系统集成演示及任务书交付对照。定稿时请将示例性数值替换为 TensorBoard、\texttt{risk\_report.json} 与整例评估脚本的实测记录，并按分层分析框架补全讨论段落，使第3章篇幅与工作量相匹配。
